% --- Archon physics formalization source begin ---
% archon:physics
% archon:covers PhyXMiniProblems/problem_phyx_mini_0844.lean
% archon:source-report reports/phyx_mini/problem_phyx_mini_0844.source.json

\chapter{Physics problem phyx\_mini\_0844}
\label{ch:PhyXMiniProblems_problem_phyx_mini_0844}

\paragraph{Problem source.}
\#\# Question

What is the electric flux through the surface?

\#\# Answer choices

- A: A: -1.2N m\textasciicircum{}2/C

- B: B: -10N m\textasciicircum{}2/C

- C: C: -2N m\textasciicircum{}2/C

- D: D: -2.3N m\textasciicircum{}2/C

\#\# Figure caption (auxiliary; use the image as primary evidence)

The image depicts a rectangular plane inclined at an angle of 30 degrees. Several magenta arrows labeled \textbackslash{}(\textbackslash{}vec\{E\}\textbackslash{}) pass through the plane, indicating an electric field with a magnitude of 180 N/C. The plane's dimensions are 15 cm by 15 cm. A black arrow labeled \textbackslash{}(\textbackslash{}hat\{n\}\textbackslash{}) is perpendicular to the plane, representing the normal vector. The arrangement and labels suggest a context related to electric flux and fields.

\#\# Dataset metadata

Electrostatics; Physical Model Grounding Reasoning, Spatial Relation Reasoning

\paragraph{Recorded answer/context.}
D: D: -2.3N m\textasciicircum{}2/C

\paragraph{Figure/image path.}
/root/proposal\_for\_physic/science-mango/phyx\_mini\_run/phyx\_data/test\_image/844.png

\paragraph{Formalization target.}
create a compiling Lean file with sorry bodies at `PhyXMiniProblems/problem\_phyx\_mini\_0844.lean`. The Lean declarations must preserve the physical quantities, dimensions or dimensional roles, figure labels, governing-law hypotheses, and final relation expressed by this problem.
Use Mathlib/Physlib names found through LeanExplore where available. If a physics API is missing, introduce faithful local abstractions rather than scalar placeholder aliases.

\begin{theorem}[Physics formalization target]
\label{thm:physics:phyx_mini_0844:target}
\lean{PhyXMiniProblems.ProblemPhyXMini0844.problem_phyx_mini_0844}
\uses{def:physics:phyx-mini-0844:phyxminiproblems-problemphyxmini0844-electricfluxinnewtonsquaremeterspercoulomb, def:physics:phyx-mini-0844:phyxminiproblems-problemphyxmini0844-electricfluxsetup, def:physics:phyx-mini-0844:phyxminiproblems-problemphyxmini0844-matchesproblemandprimaryfigure, def:physics:phyx-mini-0844:phyxminiproblems-problemphyxmini0844-hasphysicalconfiguration, def:physics:phyx-mini-0844:phyxminiproblems-problemphyxmini0844-satisfiesrectangularsurfacearealaw, def:physics:phyx-mini-0844:phyxminiproblems-problemphyxmini0844-satisfiesuniformelectricfieldmodel, def:physics:phyx-mini-0844:phyxminiproblems-problemphyxmini0844-satisfiesprimaryfiguregeometry, def:physics:phyx-mini-0844:phyxminiproblems-problemphyxmini0844-satisfiesuniformplanarelectricfluxlaw, def:physics:phyx-mini-0844:phyxminiproblems-problemphyxmini0844-recordeddatasetanswer, def:physics:phyx-mini-0844:phyxminiproblems-problemphyxmini0844-matchesdisplayedfluxtoonedecimalplace, def:physics:phyx-mini-0844:phyxminiproblems-problemphyxmini0844-isuniquematchingdisplayedchoice}
The assigned autoformalize agent should translate this physics problem into Lean declarations in the covered file, with theorem and lemma proof bodies written as `by sorry`.
\end{theorem}
\begin{proof}
This is an autoformalization task, not a proof task. Produce statements that can later be proved by the physics prover without weakening the physical model.
\end{proof}

% --- Archon declaration topology begin ---
\section{Declaration topology}
\begin{definition}[electric Field Strength Dimension]
  \label{def:physics:phyx-mini-0844:phyxminiproblems-problemphyxmini0844-electricfieldstrengthdimension}
  \lean{PhyXMiniProblems.ProblemPhyXMini0844.electricFieldStrengthDimension}
  Electric-field strength has dimension M L T⁻² C⁻¹
\end{definition}
\begin{proof}
  This is the declared model definition of electric Field Strength Dimension.
\end{proof}
\begin{definition}[electric Flux Dimension]
  \label{def:physics:phyx-mini-0844:phyxminiproblems-problemphyxmini0844-electricfluxdimension}
  \lean{PhyXMiniProblems.ProblemPhyXMini0844.electricFluxDimension}
  Electric flux has dimension M L³ T⁻² C⁻¹, equivalently N m²/C
\end{definition}
\begin{proof}
  This is the declared model definition of electric Flux Dimension.
\end{proof}
\begin{definition}[Length Quantity]
  \label{def:physics:phyx-mini-0844:phyxminiproblems-problemphyxmini0844-lengthquantity}
  \lean{PhyXMiniProblems.ProblemPhyXMini0844.LengthQuantity}
  A nonnegative, unit-independent physical length
\end{definition}
\begin{proof}
  This is the declared model definition of Length Quantity.
\end{proof}
\begin{definition}[Area Quantity]
  \label{def:physics:phyx-mini-0844:phyxminiproblems-problemphyxmini0844-areaquantity}
  \lean{PhyXMiniProblems.ProblemPhyXMini0844.AreaQuantity}
  A nonnegative, unit-independent physical area
\end{definition}
\begin{proof}
  This is the declared model definition of Area Quantity.
\end{proof}
\begin{definition}[Electric Field Strength Quantity]
  \label{def:physics:phyx-mini-0844:phyxminiproblems-problemphyxmini0844-electricfieldstrengthquantity}
  \lean{PhyXMiniProblems.ProblemPhyXMini0844.ElectricFieldStrengthQuantity}
  \uses{def:physics:phyx-mini-0844:phyxminiproblems-problemphyxmini0844-electricfieldstrengthdimension}
  A nonnegative, unit-independent electric-field magnitude
\end{definition}
\begin{proof}
  This is the declared model definition of Electric Field Strength Quantity.
\end{proof}
\begin{definition}[Electric Flux Quantity]
  \label{def:physics:phyx-mini-0844:phyxminiproblems-problemphyxmini0844-electricfluxquantity}
  \lean{PhyXMiniProblems.ProblemPhyXMini0844.ElectricFluxQuantity}
  \uses{def:physics:phyx-mini-0844:phyxminiproblems-problemphyxmini0844-electricfluxdimension}
  A signed, unit-independent electric flux
\end{definition}
\begin{proof}
  This is the declared model definition of Electric Flux Quantity.
\end{proof}
\begin{definition}[Direction3]
  \label{def:physics:phyx-mini-0844:phyxminiproblems-problemphyxmini0844-direction3}
  \lean{PhyXMiniProblems.ProblemPhyXMini0844.Direction3}
  Three-dimensional Euclidean directions used by the surface diagram
\end{definition}
\begin{proof}
  This is the declared model definition of Direction3.
\end{proof}
\begin{definition}[length Readout]
  \label{def:physics:phyx-mini-0844:phyxminiproblems-problemphyxmini0844-lengthreadout}
  \lean{PhyXMiniProblems.ProblemPhyXMini0844.lengthReadout}
  \uses{def:physics:phyx-mini-0844:phyxminiproblems-problemphyxmini0844-lengthquantity}
  Read a physical length in a selected named length unit
\end{definition}
\begin{proof}
  This is the declared model definition of length Readout.
\end{proof}
\begin{definition}[area Readout]
  \label{def:physics:phyx-mini-0844:phyxminiproblems-problemphyxmini0844-areareadout}
  \lean{PhyXMiniProblems.ProblemPhyXMini0844.areaReadout}
  \uses{def:physics:phyx-mini-0844:phyxminiproblems-problemphyxmini0844-areaquantity}
  Read a physical area in the square of a selected named length unit
\end{definition}
\begin{proof}
  This is the declared model definition of area Readout.
\end{proof}
\begin{definition}[length In Meters]
  \label{def:physics:phyx-mini-0844:phyxminiproblems-problemphyxmini0844-lengthinmeters}
  \lean{PhyXMiniProblems.ProblemPhyXMini0844.lengthInMeters}
  \uses{def:physics:phyx-mini-0844:phyxminiproblems-problemphyxmini0844-lengthquantity, def:physics:phyx-mini-0844:phyxminiproblems-problemphyxmini0844-lengthreadout}
  Metre readout of a physical length
\end{definition}
\begin{proof}
  This is the declared model definition of length In Meters.
\end{proof}
\begin{definition}[length In Centimeters]
  \label{def:physics:phyx-mini-0844:phyxminiproblems-problemphyxmini0844-lengthincentimeters}
  \lean{PhyXMiniProblems.ProblemPhyXMini0844.lengthInCentimeters}
  \uses{def:physics:phyx-mini-0844:phyxminiproblems-problemphyxmini0844-lengthquantity, def:physics:phyx-mini-0844:phyxminiproblems-problemphyxmini0844-lengthreadout}
  Centimetre readout used by both side labels in the raster
\end{definition}
\begin{proof}
  This is the declared model definition of length In Centimeters.
\end{proof}
\begin{definition}[area In Square Meters]
  \label{def:physics:phyx-mini-0844:phyxminiproblems-problemphyxmini0844-areainsquaremeters}
  \lean{PhyXMiniProblems.ProblemPhyXMini0844.areaInSquareMeters}
  \uses{def:physics:phyx-mini-0844:phyxminiproblems-problemphyxmini0844-areaquantity, def:physics:phyx-mini-0844:phyxminiproblems-problemphyxmini0844-areareadout}
  Square-metre readout of a physical area
\end{definition}
\begin{proof}
  This is the declared model definition of area In Square Meters.
\end{proof}
\begin{definition}[electric Field Strength In Newtons Per Coulomb]
  \label{def:physics:phyx-mini-0844:phyxminiproblems-problemphyxmini0844-electricfieldstrengthinnewtonspercoulomb}
  \lean{PhyXMiniProblems.ProblemPhyXMini0844.electricFieldStrengthInNewtonsPerCoulomb}
  \uses{def:physics:phyx-mini-0844:phyxminiproblems-problemphyxmini0844-electricfieldstrengthquantity}
  Coherent-SI readout of electric-field strength, in newtons/coulomb
\end{definition}
\begin{proof}
  This is the declared model definition of electric Field Strength In Newtons Per Coulomb.
\end{proof}
\begin{definition}[electric Flux In Newton Square Meters Per Coulomb]
  \label{def:physics:phyx-mini-0844:phyxminiproblems-problemphyxmini0844-electricfluxinnewtonsquaremeterspercoulomb}
  \lean{PhyXMiniProblems.ProblemPhyXMini0844.electricFluxInNewtonSquareMetersPerCoulomb}
  \uses{def:physics:phyx-mini-0844:phyxminiproblems-problemphyxmini0844-electricfluxquantity}
  Coherent-SI readout of signed electric flux, in newton-square-metres/coulomb
\end{definition}
\begin{proof}
  This is the declared model definition of electric Flux In Newton Square Meters Per Coulomb.
\end{proof}
\begin{definition}[degrees To Radians]
  \label{def:physics:phyx-mini-0844:phyxminiproblems-problemphyxmini0844-degreestoradians}
  \lean{PhyXMiniProblems.ProblemPhyXMini0844.degreesToRadians}
  Convert a dimensionless degree readout to radians
\end{definition}
\begin{proof}
  This is the declared model definition of degrees To Radians.
\end{proof}
\begin{definition}[Surface Shape]
  \label{def:physics:phyx-mini-0844:phyxminiproblems-problemphyxmini0844-surfaceshape}
  \lean{PhyXMiniProblems.ProblemPhyXMini0844.SurfaceShape}
  Shape categories relevant to the pictured surface
\end{definition}
\begin{proof}
  This is the declared model definition of Surface Shape.
\end{proof}
\begin{definition}[Surface Crossing Sense]
  \label{def:physics:phyx-mini-0844:phyxminiproblems-problemphyxmini0844-surfacecrossingsense}
  \lean{PhyXMiniProblems.ProblemPhyXMini0844.SurfaceCrossingSense}
  Which way the electric-field arrows cross relative to the displayed normal
\end{definition}
\begin{proof}
  This is the declared model definition of Surface Crossing Sense.
\end{proof}
\begin{definition}[Figure Vector Label]
  \label{def:physics:phyx-mini-0844:phyxminiproblems-problemphyxmini0844-figurevectorlabel}
  \lean{PhyXMiniProblems.ProblemPhyXMini0844.FigureVectorLabel}
  The two vector symbols printed in the figure
\end{definition}
\begin{proof}
  This is the declared model definition of Figure Vector Label.
\end{proof}
\begin{definition}[Figure Feature]
  \label{def:physics:phyx-mini-0844:phyxminiproblems-problemphyxmini0844-figurefeature}
  \lean{PhyXMiniProblems.ProblemPhyXMini0844.FigureFeature}
  Named visual features in phyx\_data/test\_image/844.png
\end{definition}
\begin{proof}
  This is the declared model definition of Figure Feature.
\end{proof}
\begin{definition}[Electric Flux Figure]
  \label{def:physics:phyx-mini-0844:phyxminiproblems-problemphyxmini0844-electricfluxfigure}
  \lean{PhyXMiniProblems.ProblemPhyXMini0844.ElectricFluxFigure}
  \uses{def:physics:phyx-mini-0844:phyxminiproblems-problemphyxmini0844-lengthquantity, def:physics:phyx-mini-0844:phyxminiproblems-problemphyxmini0844-electricfieldstrengthquantity, def:physics:phyx-mini-0844:phyxminiproblems-problemphyxmini0844-surfaceshape, def:physics:phyx-mini-0844:phyxminiproblems-problemphyxmini0844-surfacecrossingsense, def:physics:phyx-mini-0844:phyxminiproblems-problemphyxmini0844-figurevectorlabel, def:physics:phyx-mini-0844:phyxminiproblems-problemphyxmini0844-figurefeature}
  Typed data transcribed from the supplied raster. Dimension labels and the field-strength label are physical quantities; their scalar readings are stated separately in MatchesProblemAndPrimaryFigure
\end{definition}
\begin{proof}
  This is the declared model definition of Electric Flux Figure.
\end{proof}
\begin{definition}[Electric Flux Setup]
  \label{def:physics:phyx-mini-0844:phyxminiproblems-problemphyxmini0844-electricfluxsetup}
  \lean{PhyXMiniProblems.ProblemPhyXMini0844.ElectricFluxSetup}
  \uses{def:physics:phyx-mini-0844:phyxminiproblems-problemphyxmini0844-lengthquantity, def:physics:phyx-mini-0844:phyxminiproblems-problemphyxmini0844-areaquantity, def:physics:phyx-mini-0844:phyxminiproblems-problemphyxmini0844-electricfieldstrengthquantity, def:physics:phyx-mini-0844:phyxminiproblems-problemphyxmini0844-electricfluxquantity, def:physics:phyx-mini-0844:phyxminiproblems-problemphyxmini0844-direction3, def:physics:phyx-mini-0844:phyxminiproblems-problemphyxmini0844-electricfluxfigure}
  Independent physical quantities for the experiment. The requested flux is not defined from the field strength, geometry, or an answer choice; it is related to them only by the governing flux law below
\end{definition}
\begin{proof}
  This is the declared model definition of Electric Flux Setup.
\end{proof}
\begin{definition}[Matches Problem And Primary Figure]
  \label{def:physics:phyx-mini-0844:phyxminiproblems-problemphyxmini0844-matchesproblemandprimaryfigure}
  \lean{PhyXMiniProblems.ProblemPhyXMini0844.MatchesProblemAndPrimaryFigure}
  \uses{def:physics:phyx-mini-0844:phyxminiproblems-problemphyxmini0844-lengthinmeters, def:physics:phyx-mini-0844:phyxminiproblems-problemphyxmini0844-lengthincentimeters, def:physics:phyx-mini-0844:phyxminiproblems-problemphyxmini0844-electricfieldstrengthinnewtonspercoulomb, def:physics:phyx-mini-0844:phyxminiproblems-problemphyxmini0844-degreestoradians, def:physics:phyx-mini-0844:phyxminiproblems-problemphyxmini0844-electricfluxsetup}
  Problem-statement and primary-raster evidence. The raster itself, rather than the auxiliary prose, fixes the 30° as the acute angle between the field arrow and the plane. No electric-flux value or answer choice occurs here
\end{definition}
\begin{proof}
  This is the declared model definition of Matches Problem And Primary Figure.
\end{proof}
\begin{definition}[Has Physical Configuration]
  \label{def:physics:phyx-mini-0844:phyxminiproblems-problemphyxmini0844-hasphysicalconfiguration}
  \lean{PhyXMiniProblems.ProblemPhyXMini0844.HasPhysicalConfiguration}
  \uses{def:physics:phyx-mini-0844:phyxminiproblems-problemphyxmini0844-lengthinmeters, def:physics:phyx-mini-0844:phyxminiproblems-problemphyxmini0844-areainsquaremeters, def:physics:phyx-mini-0844:phyxminiproblems-problemphyxmini0844-electricfieldstrengthinnewtonspercoulomb, def:physics:phyx-mini-0844:phyxminiproblems-problemphyxmini0844-electricfluxsetup}
  Positivity, non-vacuity, and unit-direction conditions for the setup
\end{definition}
\begin{proof}
  This is the declared model definition of Has Physical Configuration.
\end{proof}
\begin{definition}[Satisfies Rectangular Surface Area Law]
  \label{def:physics:phyx-mini-0844:phyxminiproblems-problemphyxmini0844-satisfiesrectangularsurfacearealaw}
  \lean{PhyXMiniProblems.ProblemPhyXMini0844.SatisfiesRectangularSurfaceAreaLaw}
  \uses{def:physics:phyx-mini-0844:phyxminiproblems-problemphyxmini0844-lengthreadout, def:physics:phyx-mini-0844:phyxminiproblems-problemphyxmini0844-areareadout, def:physics:phyx-mini-0844:phyxminiproblems-problemphyxmini0844-electricfluxsetup}
  The area of a rectangle is the product of its side lengths. Stating the law for every selected length unit keeps the dimensionful area tied to the two dimensionful side lengths without baking in the requested flux
\end{definition}
\begin{proof}
  This is the declared model definition of Satisfies Rectangular Surface Area Law.
\end{proof}
\begin{definition}[Satisfies Uniform Electric Field Model]
  \label{def:physics:phyx-mini-0844:phyxminiproblems-problemphyxmini0844-satisfiesuniformelectricfieldmodel}
  \lean{PhyXMiniProblems.ProblemPhyXMini0844.SatisfiesUniformElectricFieldModel}
  \uses{def:physics:phyx-mini-0844:phyxminiproblems-problemphyxmini0844-electricfieldstrengthinnewtonspercoulomb, def:physics:phyx-mini-0844:phyxminiproblems-problemphyxmini0844-electricfluxsetup}
  Uniform-field idealization. Physlib's spacetime-dependent electric field is calibrated in coherent SI units on the surface by the independent physical field-strength quantity and a unit direction. This law does not mention electric flux
\end{definition}
\begin{proof}
  This is the declared model definition of Satisfies Uniform Electric Field Model.
\end{proof}
\begin{definition}[Satisfies Primary Figure Geometry]
  \label{def:physics:phyx-mini-0844:phyxminiproblems-problemphyxmini0844-satisfiesprimaryfiguregeometry}
  \lean{PhyXMiniProblems.ProblemPhyXMini0844.SatisfiesPrimaryFigureGeometry}
  \uses{def:physics:phyx-mini-0844:phyxminiproblems-problemphyxmini0844-electricfluxsetup}
  Geometry inferred directly from the primary raster. The reference direction lies in the surface and is perpendicular to the outward normal. The electric field is 30° from that in-surface direction and, because it crosses against the normal, is 90° + 30° = 120° from the outward normal
\end{definition}
\begin{proof}
  This is the declared model definition of Satisfies Primary Figure Geometry.
\end{proof}
\begin{definition}[Satisfies Uniform Planar Electric Flux Law]
  \label{def:physics:phyx-mini-0844:phyxminiproblems-problemphyxmini0844-satisfiesuniformplanarelectricfluxlaw}
  \lean{PhyXMiniProblems.ProblemPhyXMini0844.SatisfiesUniformPlanarElectricFluxLaw}
  \uses{def:physics:phyx-mini-0844:phyxminiproblems-problemphyxmini0844-areainsquaremeters, def:physics:phyx-mini-0844:phyxminiproblems-problemphyxmini0844-electricfluxinnewtonsquaremeterspercoulomb, def:physics:phyx-mini-0844:phyxminiproblems-problemphyxmini0844-electricfluxsetup}
  For a uniform electric field over a planar surface, signed electric flux is A (E · n̂). The field model above ensures that the value at the selected reference point is the common value everywhere on the surface. This is the general governing law and contains no numerical flux conclusion
\end{definition}
\begin{proof}
  This is the declared model definition of Satisfies Uniform Planar Electric Flux Law.
\end{proof}
\begin{lemma}[electric Flux exact trigonometric]
  \label{lem:physics:phyx-mini-0844:phyxminiproblems-problemphyxmini0844-electricflux-exact-trigonometric}
  \lean{PhyXMiniProblems.ProblemPhyXMini0844.electricFlux_exact_trigonometric}
  \uses{def:physics:phyx-mini-0844:phyxminiproblems-problemphyxmini0844-electricfluxinnewtonsquaremeterspercoulomb, def:physics:phyx-mini-0844:phyxminiproblems-problemphyxmini0844-electricfluxsetup, def:physics:phyx-mini-0844:phyxminiproblems-problemphyxmini0844-matchesproblemandprimaryfigure, def:physics:phyx-mini-0844:phyxminiproblems-problemphyxmini0844-hasphysicalconfiguration, def:physics:phyx-mini-0844:phyxminiproblems-problemphyxmini0844-satisfiesrectangularsurfacearealaw, def:physics:phyx-mini-0844:phyxminiproblems-problemphyxmini0844-satisfiesuniformelectricfieldmodel, def:physics:phyx-mini-0844:phyxminiproblems-problemphyxmini0844-satisfiesprimaryfiguregeometry, def:physics:phyx-mini-0844:phyxminiproblems-problemphyxmini0844-satisfiesuniformplanarelectricfluxlaw}
  Combining the rectangular area, uniform-field model, raster geometry, and planar flux law yields the exact trigonometric expression. This is a derived conclusion, not a field of any premise structure
\end{lemma}
\begin{proof}
  The conclusion follows from the governing relations and figure readouts named in the statement of electric Flux exact trigonometric.
\end{proof}
\begin{lemma}[electric Flux exact]
  \label{lem:physics:phyx-mini-0844:phyxminiproblems-problemphyxmini0844-electricflux-exact}
  \lean{PhyXMiniProblems.ProblemPhyXMini0844.electricFlux_exact}
  \uses{def:physics:phyx-mini-0844:phyxminiproblems-problemphyxmini0844-electricfluxinnewtonsquaremeterspercoulomb, def:physics:phyx-mini-0844:phyxminiproblems-problemphyxmini0844-electricfluxsetup, def:physics:phyx-mini-0844:phyxminiproblems-problemphyxmini0844-matchesproblemandprimaryfigure, def:physics:phyx-mini-0844:phyxminiproblems-problemphyxmini0844-hasphysicalconfiguration, def:physics:phyx-mini-0844:phyxminiproblems-problemphyxmini0844-satisfiesrectangularsurfacearealaw, def:physics:phyx-mini-0844:phyxminiproblems-problemphyxmini0844-satisfiesuniformelectricfieldmodel, def:physics:phyx-mini-0844:phyxminiproblems-problemphyxmini0844-satisfiesprimaryfiguregeometry, def:physics:phyx-mini-0844:phyxminiproblems-problemphyxmini0844-satisfiesuniformplanarelectricfluxlaw}
  Since cos (120°) = -1/2, the exact physical flux is -81/40 = -2.025 N m²/C
\end{lemma}
\begin{proof}
  The conclusion follows from the governing relations and figure readouts named in the statement of electric Flux exact.
\end{proof}
\begin{definition}[Answer Choice]
  \label{def:physics:phyx-mini-0844:phyxminiproblems-problemphyxmini0844-answerchoice}
  \lean{PhyXMiniProblems.ProblemPhyXMini0844.AnswerChoice}
  Labels of the four flux choices printed in the question
\end{definition}
\begin{proof}
  This is the declared model definition of Answer Choice.
\end{proof}
\begin{definition}[displayed Flux In Newton Square Meters Per Coulomb]
  \label{def:physics:phyx-mini-0844:phyxminiproblems-problemphyxmini0844-displayedfluxinnewtonsquaremeterspercoulomb}
  \lean{PhyXMiniProblems.ProblemPhyXMini0844.displayedFluxInNewtonSquareMetersPerCoulomb}
  \uses{def:physics:phyx-mini-0844:phyxminiproblems-problemphyxmini0844-answerchoice}
  Signed flux in N m²/C printed beside each answer choice
\end{definition}
\begin{proof}
  This is the declared model definition of displayed Flux In Newton Square Meters Per Coulomb.
\end{proof}
\begin{definition}[recorded Dataset Answer]
  \label{def:physics:phyx-mini-0844:phyxminiproblems-problemphyxmini0844-recordeddatasetanswer}
  \lean{PhyXMiniProblems.ProblemPhyXMini0844.recordedDatasetAnswer}
  \uses{def:physics:phyx-mini-0844:phyxminiproblems-problemphyxmini0844-answerchoice}
  The source dataset records choice D, retained here as metadata only
\end{definition}
\begin{proof}
  This is the declared model definition of recorded Dataset Answer.
\end{proof}
\begin{definition}[Matches Displayed Flux To One Decimal Place]
  \label{def:physics:phyx-mini-0844:phyxminiproblems-problemphyxmini0844-matchesdisplayedfluxtoonedecimalplace}
  \lean{PhyXMiniProblems.ProblemPhyXMini0844.MatchesDisplayedFluxToOneDecimalPlace}
  \uses{def:physics:phyx-mini-0844:phyxminiproblems-problemphyxmini0844-electricfluxinnewtonsquaremeterspercoulomb, def:physics:phyx-mini-0844:phyxminiproblems-problemphyxmini0844-electricfluxsetup, def:physics:phyx-mini-0844:phyxminiproblems-problemphyxmini0844-answerchoice, def:physics:phyx-mini-0844:phyxminiproblems-problemphyxmini0844-displayedfluxinnewtonsquaremeterspercoulomb}
  A physical flux matches a displayed one-decimal value when it is within half of 0.1 N m²/C of that value
\end{definition}
\begin{proof}
  This is the declared model definition of Matches Displayed Flux To One Decimal Place.
\end{proof}
\begin{definition}[Is Unique Matching Displayed Choice]
  \label{def:physics:phyx-mini-0844:phyxminiproblems-problemphyxmini0844-isuniquematchingdisplayedchoice}
  \lean{PhyXMiniProblems.ProblemPhyXMini0844.IsUniqueMatchingDisplayedChoice}
  \uses{def:physics:phyx-mini-0844:phyxminiproblems-problemphyxmini0844-electricfluxsetup, def:physics:phyx-mini-0844:phyxminiproblems-problemphyxmini0844-answerchoice, def:physics:phyx-mini-0844:phyxminiproblems-problemphyxmini0844-matchesdisplayedfluxtoonedecimalplace}
  A choice is the sole displayed value matching the independently modeled flux
\end{definition}
\begin{proof}
  This is the declared model definition of Is Unique Matching Displayed Choice.
\end{proof}
% --- Archon declaration topology end ---
% --- Archon physics formalization source end ---
